% 06_appendices.tex — Apéndices A–F (material suplementario)
%   A: consultas por base  B: prompt y sensibilidad  C: riesgo de sesgo
%   D: notas sobre κ       E: lista PRISMA-ScR      F: evidencia complementaria
% Los datos aún no disponibles se definen en 00_datos_pendientes.tex.

\clearpage
\nolinenumbers
\appendix
\renewcommand*{\appendixname}{Appendix}
\addto\captionsspanish{\renewcommand*{\appendixname}{Appendix}}
\titleformat{\section}{\normalsize\bfseries}{\appendixname~\thesection.}{1em}{}
\counterwithin{table}{section}
\counterwithin{figure}{section}
\renewcommand{\thetable}{\thesection\arabic{table}}
\renewcommand{\thefigure}{\thesection\arabic{figure}}

% ─────────────────────────────────────────────────────────────────────────────
\section{Consultas booleanas por base de datos}
\label{apx:queries_multibase}
% ─────────────────────────────────────────────────────────────────────────────

Se reproducen las tres consultas booleanas, mapeadas 1-a-1 con las tres
preguntas de investigación (RQ1 modelado y sensores virtuales; RQ2 control
avanzado frente a PID; RQ3 ML/RL/híbridos y barreras de adopción). Todas
las cadenas se ejecutaron el \FechaBusqueda{} sobre las cinco bases
de datos (Scopus, Web of Science, IEEE~Xplore, SciELO, LILACS) con ventana
temporal 2015--2026 (incluyendo registros \textit{early access} con año
editorial 2026). La sintaxis por base se ajustó al motor correspondiente
conservando el mismo alcance semántico: Scopus utilizó el campo
\texttt{TITLE-ABS-KEY(\ldots)}; Web of Science empleó
\texttt{TS=(\ldots)}~(\textit{Topic}); IEEE~Xplore usó operadores
directos sobre \textit{All Metadata}; SciELO empleó el campo
\texttt{(ti,ab)} con operadores booleanos estándar; LILACS empleó
búsqueda por DeCS (Descriptores en Ciencias de la Salud) con operadores
booleanos y \texttt{AND NOT} sustituido por \texttt{NOT}.

\begin{table}[H]
\centering
\caption{Recuentos verificados por base de datos y consulta ($N=541$
en las tres bases originales; \NSciELO{} en SciELO y
\NLILACS{} en LILACS; búsqueda del \FechaBusqueda, ventana 2015--2026).}
\label{tab:queries_multibase}
\singlespacing\small
\begin{tabular}{l l p{8.2cm} r}
\toprule
\textbf{Base} & \textbf{Consulta} & \textbf{Foco (RQ)} & \textbf{$n$} \\
\midrule
\multirow{4}{*}{Scopus}
 & Q1 & pH, sensores virtuales e identificación (RQ1) & 11 \\
 & Q2 & Control avanzado y metaheurísticas vs.\ PID (RQ2) & 32 \\
 & Q3 & ML/RL/híbridos y barreras (RQ3) & 269 \\
 & & \textbf{Subtotal Scopus} & \textbf{312} \\
\midrule
\multirow{4}{*}{Web of Science}
 & Q1 & pH, sensores virtuales e identificación (RQ1) & 3 \\
 & Q2 & Control avanzado y metaheurísticas vs.\ PID (RQ2) & 18 \\
 & Q3 & ML/RL/híbridos y barreras (RQ3) & 167 \\
 & & \textbf{Subtotal WoS} & \textbf{188} \\
\midrule
\multirow{4}{*}{IEEE~Xplore}
 & Q1 & pH, sensores virtuales e identificación (RQ1) & 4 \\
 & Q2 & Control avanzado y metaheurísticas vs.\ PID (RQ2) & 10 \\
 & Q3 & ML/RL/híbridos y barreras (RQ3) & 27 \\
 & & \textbf{Subtotal IEEE} & \textbf{41} \\
\midrule
\multicolumn{3}{l}{\textbf{Total bruto original (3 bases)}} & \textbf{541} \\
\midrule
% ════════════════════════════════════════════════════════════════════════════
% NUEVO v2.0 — Filas para SciELO y LILACS
% ════════════════════════════════════════════════════════════════════════════
\multirow{2}{*}{SciELO}
 & Q4 & Adaptación combinada de Q1+Q2+Q3 a sintaxis (ti,ab) con
      operadores booleanos & \NSciELO \\
 & & \textbf{Subtotal SciELO} & \textbf{\NSciELO} \\
\midrule
\multirow{2}{*}{LILACS}
 & Q4 & Adaptación combinada a DeCS (Descriptores en Ciencias de la
      Salud) con operadores booleanos & \NLILACS \\
 & & \textbf{Subtotal LILACS} & \textbf{\NLILACS} \\
\midrule
\multicolumn{3}{l}{\textbf{Total bruto ampliado (5 bases)}} & \textbf{\NTotalAmpliado} \\
\multicolumn{3}{l}{Tras deduplicación por DOI y \textit{fuzzy match} de títulos}
 & \textbf{\NUnicosAmpliado} \\
\bottomrule
\end{tabular}
\begin{Nota}[Nota]{\linewidth}{La deduplicación combinó emparejamiento
exacto por DOI y \textit{fuzzy match} de títulos con umbral $\geq 95\,\%$
(\textit{rapidfuzz}), reduciendo falsos negativos por variaciones
tipográficas entre RIS de distintas fuentes. La tasa de solapamiento
inter-base en las tres bases originales ($34.8\,\%=188/541$) es consistente
con una estrategia de búsqueda estable y con convergencia de un núcleo
común de literatura entre las tres bases. Las consultas SciELO y LILACS
se diseñaron como una única consulta combinada Q4 por base, adaptando los
elementos semánticos comunes de Q1+Q2+Q3 a la sintaxis local (operadores
booleanos, campos disponibles).}
\end{Nota}
\end{table}

% ════════════════════════════════════════════════════════════════════════════
% NUEVO v2.0 — Tabla A2 con sintaxis específica SciELO y LILACS
% ════════════════════════════════════════════════════════════════════════════
\newpage

\begin{table}[H]
\centering
\caption{Cadenas específicas para SciELO y LILACS con sintaxis local
(consulta combinada Q4; elementos semánticos comunes de Q1+Q2+Q3).}
\label{tab:scielo_lilacs_syntax}
\singlespacing\scriptsize
\begin{tabular}{p{2.3cm} p{12.2cm}}
\toprule
\textbf{Base} & \textbf{Cadena (sintaxis local)} \\
\midrule
SciELO
  & \texttt{(ti:("pH control" OR "pH monitoring" OR "pH prediction"
     OR "pH regulation" OR "acidity control" OR "alkalinity control"
     OR "lime dosage" OR "calcium hydroxide" OR "milk of lime") AND
     ab:("sugar mill" OR "sugarcane" OR "cane juice" OR "sugar juice"
     OR "sugar industry" OR "juice clarification" OR "liming process"
     OR "defecation" OR "sulfitation" OR "sugar beet" OR "sugar factory"
     OR "saccharum") AND (ti:("soft sensor" OR "virtual sensor" OR
     "identification" OR "neural network" OR "LSTM" OR "predictive model"
     OR "model predictive control" OR "NMPC" OR "fuzzy logic" OR
     "machine learning" OR "deep learning" OR "reinforcement learning"
     OR "digital twin") OR ab:("soft sensor" OR "virtual sensor" OR
     "identification" OR "neural network" OR "LSTM" OR "predictive model"
     OR "model predictive control" OR "NMPC" OR "fuzzy logic" OR
     "machine learning" OR "deep learning" OR "reinforcement learning"
     OR "digital twin"))) NOT (ti:("wastewater" OR "soil pH" OR
     "fermentation broth" OR "bioethanol" OR "pharmaceutical" OR
     "bioreactor" OR "wine" OR "beer" OR "dairy") OR ab:("wastewater"
     OR "soil pH" OR "fermentation broth" OR "bioethanol" OR
     "pharmaceutical" OR "bioreactor" OR "wine" OR "beer" OR "dairy"))} \\
\midrule
LILACS (DeCS)
  & \texttt{("Controle de pH" OR "monitoramento de pH" OR "previsão de pH"
     OR "regulação de pH" OR "controle de acidez" OR "dosagem de cal"
     OR "hidróxido de cálcio" OR "leite de cal" OR "pH control" OR
     "pH monitoring" OR "pH prediction" OR "pH regulation" OR
     "acidity control" OR "lime dosage" OR "calcium hydroxide" OR
     "milk of lime") AND ("usina de açúcar" OR "cana-de-açúcar" OR
     "caldo de cana" OR "suco de cana" OR "indústria açucareira" OR
     "clarificação de caldo" OR "calagem" OR "defecação" OR
     "sulfitação" OR "beterraba açucareira" OR "sugar mill" OR
     "sugarcane" OR "cane juice" OR "sugar juice" OR "sugar industry"
     OR "juice clarification" OR "liming process" OR "defecation" OR
     "sulfitation" OR "sugar beet" OR "sugar factory" OR "saccharum")
     AND ("sensor virtual" OR "sensor suave" OR "identificação de
     sistemas" OR "rede neural" OR "LSTM" OR "modelo preditivo" OR
     "controle preditivo de modelo" OR "NMPC" OR "lógica fuzzy" OR
     "aprendizado de máquina" OR "aprendizado profundo" OR
     "aprendizado por reforço" OR "gêmeo digital" OR "soft sensor"
     OR "virtual sensor" OR "identification" OR "neural network" OR
     "model predictive control" OR "fuzzy logic" OR "machine learning"
     OR "deep learning" OR "reinforcement learning" OR "digital twin")
     NOT ("águas residuais" OR "pH do solo" OR "caldo de fermentação"
     OR "bioetanol" OR "farmacêutico" OR "biorreator" OR "vinho" OR
     "cerveja" OR "laticínios" OR "wastewater" OR "soil pH" OR
     "fermentation broth" OR "bioethanol" OR "pharmaceutical" OR
     "bioreactor" OR "wine" OR "beer" OR "dairy")} \\
\bottomrule
\end{tabular}
\begin{Nota}[Nota]{\linewidth}{Elaboración propia. La consulta Q4
combina los elementos semánticos comunes de Q1, Q2 y Q3 en una sola
cadena por base, ajustándose a la sintaxis local. SciELO emplea operadores
\texttt{(ti:)} (título) y \texttt{(ab:)} (resumen) con booleanos estándar;
\texttt{AND NOT} se sustituye por \texttt{NOT}. LILACS emplea descriptores
DeCS bilingües (portugués + español) para máxima sensibilidad regional;
los términos en inglés se mantienen como refuerzo. Los recuentos se
reportan en la Tabla~\ref{tab:queries_multibase}.}
\end{Nota}
\end{table}

\newpage

\subsection*{Q1 --- Modelado y sensores virtuales de pH (RQ1)}

\begin{lstlisting}[language={},caption={Cadena Q1 (Scopus). Las cadenas
equivalentes para WoS e IEEE~Xplore se obtienen sustituyendo el operador
de campo (\texttt{TS=} en WoS; sin prefijo en IEEE).},label={lst:q1}]
TITLE-ABS-KEY(
  ("pH control" OR "pH monitoring" OR "pH prediction"
   OR "pH optimization" OR "pH regulation" OR "pH measurement"
   OR "pH adjustment" OR "pH correction" OR "acidity control"
   OR "alkalinity control" OR "juice acidity"
   OR "lime dosage" OR "lime addition"
   OR "calcium hydroxide" OR "milk of lime")
  AND
  ("sugar mill" OR "sugarcane" OR "cane juice"
   OR "sugar juice" OR "sugar industry"
   OR "juice clarification" OR "liming process"
   OR "defecation" OR "sulfitation"
   OR "sugar beet" OR "beet sugar"
   OR "sugar factory" OR "saccharum")
  AND
  ("soft sensor" OR "virtual sensor" OR "inferential sensor"
   OR "identification" OR "data-driven model" OR "data-driven"
   OR "online estimation" OR "real-time estimation"
   OR "neural network" OR "LSTM"
   OR "predictive model" OR "regression model")
)
AND NOT TITLE-ABS-KEY(
  "wastewater" OR "water treatment" OR "soil pH" OR "blood pH"
  OR "fermentation broth" OR "bioethanol"
  OR "pharmaceutical" OR "bioreactor"
  OR "wine" OR "beer" OR "dairy" OR "groundwater"
)
\end{lstlisting}

\subsection*{Q2 --- Control avanzado (NMPC/Fuzzy/metaheurísticas) vs.\ PID (RQ2)}

\begin{lstlisting}[language={},caption={Cadena Q2 (Scopus).},label={lst:q2}]
TITLE-ABS-KEY(
  ("model predictive control" OR "NMPC" OR "nonlinear MPC"
   OR "fuzzy logic" OR "takagi-sugeno" OR "neuro-fuzzy"
   OR "adaptive control" OR "intelligent control"
   OR "PID tuning" OR "optimal control"
   OR "pattern search" OR "particle swarm"
   OR "genetic algorithm" OR "differential evolution"
   OR "bayesian optimization")
  AND "pH"
  AND ("sugar mill" OR "sugarcane" OR "cane juice"
       OR "sugar juice" OR "juice clarification"
       OR "liming process" OR "defecation" OR "sulfitation"
       OR "sugar industry" OR "sugar beet" OR "sugar factory")
)
AND NOT TITLE-ABS-KEY(
  "wastewater" OR "bioethanol" OR "fermentation broth"
  OR "pharmaceutical" OR "soil" OR "blood"
  OR "wine" OR "dairy" OR "desalination"
  OR "anaerobic digestion" OR "bioreactor"
)
\end{lstlisting}

\subsection*{Q3 --- ML/RL/híbridos y barreras de adopción (RQ3)}

\begin{lstlisting}[language={},caption={Cadena Q3 (Scopus).},label={lst:q3}]
TITLE-ABS-KEY(
  ("machine learning" OR "deep learning" OR "neural network"
   OR "artificial intelligence" OR "random forest"
   OR "gradient boosting" OR "XGBoost" OR "lightGBM"
   OR "support vector"
   OR "reinforcement learning" OR "deep reinforcement learning"
   OR "DDPG" OR "PPO" OR "Q-learning"
   OR "hybrid model" OR "grey-box model" OR "physics-informed"
   OR "digital twin" OR "transfer learning")
  AND "pH"
  AND ("sugar" OR "sugarcane" OR "cane juice"
       OR "sugar juice" OR "clarification"
       OR "liming" OR "sugar mill" OR "sugar industry"
       OR "sugar beet" OR "juice purification"
       OR "sugar factory" OR "defecation")
)
AND NOT TITLE-ABS-KEY(
  "wastewater" OR "bioethanol" OR "pharmaceutical"
  OR "soil" OR "blood" OR "bioreactor" OR "desalination"
  OR "wine" OR "beer" OR "fermentation broth"
)
\end{lstlisting}

% ─────────────────────────────────────────────────────────────────────────────
\section{\textit{Prompt} del chequeo automático complementario y análisis de
sensibilidad}
\label{apx:prompt}
% ─────────────────────────────────────────────────────────────────────────────

\noindent\textbf{Función y advertencias.} El \textit{prompt} que sigue se
empleó \emph{exclusivamente} como chequeo automático complementario
(Etapa~5, \ref{sub:seleccion}); no es una etapa de cribado ni sustituye el
juicio humano, cuya autoridad reside en el doble cribado de la Etapa~4. Se
reproduce tal como lo ejecuta el \textit{script} \texttt{04\_filtrado\_ollama.py}
(modelo servido por Ollama con $T=0$; salida JSON). Conviene precisar tres
hechos para su correcta interpretación:
\begin{enumerate}[leftmargin=*,noitemsep]
\item Los seis ejemplos son \textbf{sintéticos}: no corresponden a
publicaciones reales, y sus cifras (p.~ej., ``MAE$=0.04$'') son
ilustrativas.
\item El \textbf{Ejemplo~3} reproduce el título, un resumen parafraseado y
las palabras clave de uno de los estudios incluidos
\parencite{palaciomorales2021optimiza}. Es una \textbf{fuga de información}:
el modelo recibió como ejemplo positivo el estudio que luego clasificó
como positivo. Su efecto se evaluó re-ejecutando el \textit{pipeline} sin
la fuga (Tabla~\ref{tab:sens_prompt}).
\item El criterio de cuartil SJR \textbf{no forma parte del
\textit{prompt}}; se aplica antes (Etapa~3).
\end{enumerate}

\begin{lstlisting}[language={},
caption={\textit{Prompt} de producción (seis ejemplos sintéticos) para
clasificación INCLUIR/EXCLUIR. El Ejemplo~3 contiene la fuga descrita en
el texto.},
label={lst:prompt}]
Actua como revisor experto en una revision sistematica
(RSL) sobre APRENDIZAJE AUTOMATICO aplicado al CONTROL DE
pH EN LA INDUSTRIA AZUCARERA (clarificacion, encalado,
sulfitacion, purificacion del jugo de cana y remolacha).

=== EJEMPLOS DE REFERENCIA (six-shot) ===

EJEMPLO 1 - INCLUIR (sintetico: ANN para prediccion de pH)
Titulo: "Artificial neural network for real-time pH
         prediction in sugarcane juice clarification: lime
         dosage optimization in a Brazilian sugar mill"
Resumen: "A multilayer perceptron trained on 24 months of
         DCS data predicts clarified juice pH ... MAE=0.04
         pH units. Reduces lime overconsumption by 18%."
Decision correcta: INCLUIR   Confianza: 97

EJEMPLO 2 - INCLUIR (sintetico: LSTM + soft sensor)
Titulo: "LSTM-based soft sensor for online pH estimation in
         sugar cane juice purification without physical pH
         probe"
Resumen: "An LSTM model estimates juice pH from conductivity,
         temperature, turbidity, Brix and lime flow ...
         R2=0.96, RMSE=0.05 ..."
Decision correcta: INCLUIR   Confianza: 96

EJEMPLO 3 - INCLUIR (FUGA: reproduce un estudio incluido)
Titulo: "Optimization based on pattern search algorithm
         applied to pH non-linear control: Application to
         alkalinization process of sugar juice"
Resumen: "A Takagi-Sugeno fuzzy inference model combined with
         a non-linear model predictive controller (NMPC)
         tuned via pattern search optimization (PSM) controls
         pH in the alkalinization of sugar juice ..."
Keywords: pH control; alkalinization; sugar juice; NMPC;
         fuzzy; pattern search; nonlinear
Decision correcta: INCLUIR   Confianza: 95

EJEMPLO 4 - EXCLUIR (sintetico: pH en bioetanol de cana)
Titulo: "Machine learning control of fermentation pH for
         bioethanol production from sugarcane bagasse"
Decision correcta: EXCLUIR   Confianza: 96

EJEMPLO 5 - EXCLUIR (sintetico: pH en cristalizacion)
Titulo: "Neural network prediction of massecuite pH and
         crystal size in vacuum pan"
Decision correcta: EXCLUIR   Confianza: 91

EJEMPLO 6 - EXCLUIR (sintetico: aguas residuales)
Titulo: "Deep learning for pH prediction in sugar factory
         wastewater treatment"
Decision correcta: EXCLUIR   Confianza: 95

=== FIN DE EJEMPLOS ===

INCLUIR si el articulo cumple TODOS:
  A) Aplica tecnica de ML/IA (machine learning, redes
     neuronales, soft sensors, reinforcement learning, MPC,
     LSTM, CNN, fuzzy logic, gemelos digitales, etc.).
  B) La aplicacion es CONTROL, PREDICCION, ESTIMACION,
     MONITOREO u OPTIMIZACION del pH (o acidez, alcalinidad,
     dosis de cal, dosis de SO2) mediante ML/IA.
  C) El proceso objetivo PRINCIPAL es la INDUSTRIA AZUCARERA:
     clarificacion, encalado, defecacion, sulfitacion,
     carbonatacion o purificacion del jugo de cana o remolacha.

EXCLUIR si se cumple CUALQUIERA:
  X1) pH en procesos NO azucareros.   X2) pH en suelo, agua
  subterranea, rios o ambiente.   X3) Procesos azucareros
  distintos a clarificacion/encalado/sulfitacion.   X4)
  Contexto multi-industria alimentaria.   X5) pH solo como
  variable de entrada.   X6) Tejidos biologicos o farmacia.
  X7) Efluentes de la fabrica azucarera.   X8) Bioetanol o
  biocombustibles de cana.

Responde UNICAMENTE con este JSON exacto:
{"decision": "INCLUIR" o "EXCLUIR",
 "confianza": numero entre 0 y 100,
 "justificacion": "texto conciso, maximo 2 oraciones"}

=== ARTICULO A EVALUAR ===
Titulo: {titulo}   Resumen: {resumen}   Palabras clave: {keywords}
\end{lstlisting}

\noindent\textbf{Parámetros de ejecución y reproducibilidad.} Modelo
\texttt{llama3:latest} (Llama~3 de 8B parámetros, cuantización Q4\_0)
servido localmente con Ollama; resumen SHA-256 del modelo:
\texttt{365c0bd3c000a25d28ddbf732fe1c6add414de7275464c4e4d1c3b5fcb5d8ad1}
(fecha de modificación local: 23 de marzo de 2026, anterior a la corrida de producción);
versión del servidor en la re-ejecución: 0.24.0; temperatura $T=0$; no se
fijó semilla. Umbral de derivación: confianza $<55\,\%$ (tratado como
EXCLUIR). En la corrida de producción sobre los 230~registros, 201 se
resolvieron por el filtro léxico local sin invocar al modelo (nivel
\texttt{LOCAL}), 28 obtuvieron una decisión de alta confianza (nivel
\texttt{AUTO}, $92$--$98\,\%$; 2~INCLUIR y 26~EXCLUIR) y 1 cayó bajo el
umbral (nivel \texttt{INCIERTO}; una entrada de actas mal etiquetada).

\subsection*{Análisis de sensibilidad del \textit{prompt}}

Los 29~registros que llegaron al modelo se re-procesaron con el mismo
modelo, $T=0$, con (A)~el \textit{prompt} original y (B)~un \textit{prompt}
corregido en el que el Ejemplo~3 se sustituye por un positivo sintético de
otra familia metodológica (control neuro-difuso con programación de
ganancia). Las 26 decisiones EXCLUIR restantes no cambiaron. La
Tabla~\ref{tab:sens_prompt} resume los registros con decisión relevante.

\begin{table}[H]
\centering
\caption{Decisión del modelo por registro bajo tres ejecuciones
(producción; re-ejecución con \textit{prompt} original A; re-ejecución con
\textit{prompt} sin fuga B). Entre paréntesis, confianza autoinformada.}
\label{tab:sens_prompt}
\singlespacing\small
\begin{tabular}{p{6.0cm} c c c c}
\toprule
\textbf{Registro} & \textbf{Humano} & \textbf{Producción} & \textbf{A} & \textbf{B} \\
\midrule
\cite{palaciomorales2021optimiza} & INCLUIR & INCLUIR (95) & INCLUIR (95) & INCLUIR (98) \\
\cite{murugesan2020computeraided} & INCLUIR & INCLUIR (98) & INCLUIR (98) & INCLUIR (98) \\
Clarificación de jugo con H$_2$O$_2$ y redes neuronales (2017) & EXCLUIR & EXCLUIR (92) & INCLUIR (98) & INCLUIR (98) \\
\midrule
Otros 26 registros & EXCLUIR & EXCLUIR & EXCLUIR & EXCLUIR \\
\bottomrule
\end{tabular}
\begin{Nota}[Nota]{\linewidth}{Concordancia con la decisión humana final:
producción 230/230; A y B 229/230. La fuga no determinó la clasificación
del estudio de \cite{palaciomorales2021optimiza}, pero la decisión sobre el
registro de H$_2$O$_2$ cambia entre corridas con $T=0$, de modo que el
chequeo no es reproducible.}
\end{Nota}
\end{table}

% ─────────────────────────────────────────────────────────────────────────────
\section{Herramienta de evaluación del riesgo de sesgo}
\label{apx:herramienta_sesgo}
% ─────────────────────────────────────────────────────────────────────────────

Se diseñó una herramienta \textit{ad~hoc} de cinco dominios, adaptada de
los principios generales de evaluación del riesgo de sesgo en estudios
primarios (representatividad de los datos, separación entrenamiento/
validación, completitud de las métricas, línea base y nivel de validación)
al contexto de AA y control avanzado en procesos industriales, para el
cual no existe una herramienta estandarizada. Su aplicación es opcional
en PRISMA-ScR y no se usa para ponderar resultados. La
Tabla~\ref{tab:rob_dominios} define los dominios y la
Tabla~\ref{tab:sesgo} los aplica a los dos estudios incluidos, a partir
del \textbf{texto completo}.

\begin{table}[H]
\centering
\caption{Dominios de la herramienta de evaluación del riesgo de sesgo.}
\label{tab:rob_dominios}
\singlespacing\small
\begin{tabular}{p{3.6cm} p{5.0cm} p{5.0cm}}
\toprule
\textbf{Dominio} & \textbf{Pregunta guía} & \textbf{Escala} \\
\midrule
D1 --- Representatividad &
  ¿Los datos usados para identificar/entrenar el modelo provienen del
  proceso industrial real, de un banco de pruebas, o exclusivamente de
  simulación? &
  Industrial $>$ banco de pruebas $>$ simulación \\
D2 --- Separación de datos &
  ¿Se declara una separación explícita entre los datos usados para
  ajustar el modelo/controlador y los usados para validarlo (p.~ej.,
  \textit{train/test split}, validación cruzada)? &
  Declarada y adecuada / declarada pero limitada / no declarada \\
D3 --- Completitud de métricas &
  ¿Se reportan métricas cuantitativas de desempeño del control de pH
  (p.~ej., error de seguimiento, \textit{overshoot}, tiempo de
  establecimiento) con suficiente detalle para ser reproducidas? &
  $\geq 2$ métricas cuantitativas / 1 métrica o solo cualitativa /
  ninguna \\
D4 --- Comparación con línea base &
  ¿El desempeño del método propuesto se compara cuantitativamente
  contra un controlador o algoritmo de referencia (p.~ej., PID, PI, u
  otra metaheurística)? &
  Comparación cuantitativa explícita / comparación solo cualitativa /
  sin comparación \\
D5 --- Nivel de validación alcanzado &
  ¿En qué entorno se validó el controlador o modelo resultante
  (\textit{Technology Readiness Level} aproximado, escala NASA
  NPR~7123.1B)? &
  Planta industrial (TRL~$\geq$~7) / planta piloto (TRL~5--6) /
  simulación (TRL~$\leq$~4) \\
\bottomrule
\end{tabular}
\begin{Nota}[Nota]{\linewidth}{El riesgo global de cada estudio se
calificó de forma cualitativa como \textit{Bajo}, \textit{Moderado} o
\textit{Moderado--Alto} según el número de dominios en su categoría más
favorable, sin un algoritmo de puntuación numérica formal, dado el
tamaño reducido de la muestra de estudios evaluados ($n=2$) y el
carácter exploratorio de esta revisión de alcance.}
\end{Nota}
\end{table}

\begin{table}[H]
\centering
\caption{Aplicación de la herramienta de riesgo de sesgo a los dos
estudios incluidos (texto completo).}
\label{tab:sesgo}
\singlespacing\small
\begin{tabular}{p{2.6cm} p{2.4cm} p{2.4cm} p{2.4cm} p{2.4cm} p{2.4cm}}
\toprule
\textbf{Estudio} & \textbf{D1} & \textbf{D2} & \textbf{D3} &
  \textbf{D4} & \textbf{D5} \\
\midrule
\cite{palaciomorales2021optimiza} &
  Simulación (modelo genérico de neutralización) &
  No declarada &
  $\geq 2$ métricas (ICET, sobreimpulso, esfuerzo) &
  Contra otras cinco sintonías de NMPC; sin PID/PI &
  Simulación (TRL~$\approx 3$--$4$) \\
\cite{murugesan2020computeraided} &
  Datos reales de planta (3000 muestras) &
  Declarada (70/30\,\%) &
  $\geq 2$ métricas (MSE, tiempos de subida y asentamiento) &
  Contra GA y PSO; sin PID de planta &
  Controlador solo en simulación (TRL~$\approx 4$) \\
\bottomrule
\end{tabular}
\begin{Nota}[Nota]{\linewidth}{\textbf{Riesgo global:}
\cite{palaciomorales2021optimiza} --- \textit{Moderado--Alto}: D1, D2 y D5
son desfavorables (modelo de planta simulado, sin separación de datos
declarada, sin planta), D3 es favorable y D4 parcial; sin comparación con
un controlador de referencia tipo PID.
\cite{murugesan2020computeraided} --- \textit{Moderado}: D1, D2 y D3 son
favorables y D4 parcial (comparación solo con otras metaheurísticas),
pero D5 es desfavorable porque el controlador no se probó en planta y el
modelo identificado es lineal de segundo orden. Ambos comparten la
ausencia de validación en lazo cerrado en planta, coherente con la
prioridad que encabeza la agenda (\ref{sub:agenda}). Una versión
anterior de esta evaluación se basó en los resúmenes y se corrigió tras la
lectura del texto completo.}
\end{Nota}
\end{table}

% ─────────────────────────────────────────────────────────────────────────────
\section{Notas metodológicas sobre la fiabilidad del cribado
(\texorpdfstring{$\kappa$}{kappa})}
\label{apx:kappa_notas}
% ─────────────────────────────────────────────────────────────────────────────

\subsection*{1. Métrica primaria: $\kappa$ inter-humano}

La concordancia entre los Revisores~1 y~2 se calcula sobre los
230~registros (Tabla~\ref{tab:confusion_inter}); ambos cribaron el corpus
completo, por lo que no hay error de muestreo en la cobertura. Con solo
dos casos positivos, sin embargo, un único desacuerdo habría bajado
$\kappa$ de forma drástica (con un desacuerdo sobre un positivo,
$\kappa\approx0.66$), y la ausencia total de desacuerdos produce
$\kappa=1.000$ y $\text{PABAK}=1.000$ sin que ello permita estimar la
variabilidad entre revisores. Por eso la métrica se reporta junto con
$P_0$ y los intervalos de Wilson (sensibilidad y precisión: 2/2,
IC 95\,\% 34.2--100\,\%; $P_0$: 230/230, IC 95\,\% 98.4--100\,\%), y se
interpreta como estimación de la fiabilidad del proceso de exclusión,
no del acierto sobre los positivos
\parencite{feinstein1990high,byrt1993bias}.

\subsection*{2. Métrica secundaria: concordancia con el \textit{pipeline}
automático}

\begin{table}[H]
\centering
\caption{Concordancia del \textit{pipeline} automático con la decisión
humana final sobre los 230~registros Q1--Q4: corrida de producción y
re-ejecución (\textit{prompts} A y B).}
\label{tab:confusion_llm}
\singlespacing\small
\begin{tabular}{l c c c c c c}
\toprule
 & VP & FP & FN & VN & $P_0$ & $\kappa$ \\
\midrule
Producción & 2 & 0 & 0 & 228 & 1.000 & 1.000 \\
Re-ejecución A y B & 2 & 1 & 0 & 227 & 0.996 & 0.798 \\
\bottomrule
\end{tabular}
\begin{Nota}[Nota]{\linewidth}{VP: verdaderos positivos; FP: falsos
positivos; FN: falsos negativos; VN: verdaderos negativos. Producción:
sensibilidad 100\,\% (IC 95\,\% de Wilson $[34.2,100]$), precisión
100\,\% ($[34.2,100]$), $F_1=100\,\%$, PABAK $=1.000$. Re-ejecución:
sensibilidad 100\,\% ($[34.2,100]$), precisión 66.7\,\% ($[20.8,93.9]$),
$F_1=80.0\,\%$, PABAK $=0.991$.}
\end{Nota}
\end{table}

\subsection*{3. Fragilidad estadística y paradoja del $\kappa$}

Con dos observaciones positivas, la sensibilidad y la precisión provienen
de cocientes 2/2 o 2/3 y sus intervalos de Wilson al 95\,\% son amplios
($[34.2\,\%,100\,\%]$ y $[20.8\,\%,93.9\,\%]$). La diferencia entre
$\kappa=0.798$ y $\text{PABAK}=0.991$ en la re-ejecución ilustra que
$\kappa$ cae con la prevalencia mientras PABAK no, y que ninguno
sustituye a la inspección de la matriz de confusión. Un clasificador
que respondiera EXCLUIR siempre alcanzaría $P_0=0.991$ y
$\text{PABAK}=0.983$ sin detectar ningún positivo. La inestabilidad
$1.000\to0.798$ entre corridas con $T=0$ confirma que no debe presentarse
$\kappa=1.000$ como resultado definitivo.

\subsection*{4. Requisitos para una validación robusta}

(a)~Incluir un tercer revisor o estimar la concordancia en un corpus
con más positivos, para que $\kappa$ sea informativo; (b)~fijar la versión
y el resumen SHA-256 del modelo y una semilla, y repetir varias corridas
para estimar la variabilidad; (c)~sustituir el Ejemplo~3 por un positivo
sintético sin fuga; (d)~comparar varios modelos; y (e)~repetir la
evaluación en un corpus de mayor prevalencia. La reducción teórica de
carga manual ($228/230\approx 99.1\,\%$) es un beneficio potencial, no un
ahorro materializado, pues la revisión humana cubrió el 100\,\% del
corpus.

% ─────────────────────────────────────────────────────────────────────────────
\section{Lista de verificación PRISMA-ScR}
\label{apx:prisma}
% ─────────────────────────────────────────────────────────────────────────────

\singlespacing\small
\begin{longtable}{>{\raggedright\arraybackslash}p{0.7cm} >{\raggedright\arraybackslash}p{6.2cm} >{\raggedright\arraybackslash}p{7.6cm}}
\caption{Lista de verificación PRISMA-ScR \parencite{tricco2018prismascr}
y ubicación de cada ítem en el manuscrito.}\label{tab:prisma_checklist}\\
\toprule
\textbf{\#} & \textbf{Ítem} & \textbf{Ubicación} \\
\midrule
\endfirsthead
\toprule
\textbf{\#} & \textbf{Ítem} & \textbf{Ubicación} \\
\midrule
\endhead
\bottomrule
\endfoot
1 & Título & Portada: identifica el diseño como revisión de alcance \\
2 & Resumen estructurado & Resumen (párrafo único con contexto, métodos, resultados y conclusión) \\
3 & Justificación & \S\ref{sub:brecha}; \S\ref{sub:otros_dominios} \\
4 & Objetivos & \S\ref{sub:rqs} (PICOC y RQ1--RQ3) \\
5 & Protocolo y registro & \S\ref{sub:diseno} (OSF, registro retrospectivo; desviaciones declaradas) \\
6 & Criterios de elegibilidad & \S\ref{sub:criterios} \\
7 & Fuentes de información & \S\ref{sub:fuentes} \\
8 & Estrategia de búsqueda & \S\ref{sub:fuentes}; Apéndice~\ref{apx:queries_multibase} \\
9 & Selección de las fuentes de evidencia & \S\ref{sub:seleccion}; \S\ref{sub:fiabilidad_metodos} \\
10 & Proceso de extracción de datos & \S\ref{sub:extraccion} \\
11 & Elementos de datos & \S\ref{sub:extraccion}; Tabla~\ref{tab:studies} \\
12 & Evaluación crítica de las fuentes (opcional) & \S\ref{sub:extraccion}; Apéndice~\ref{apx:herramienta_sesgo} \\
13 & Síntesis de resultados & \S\ref{sub:extraccion}; \S\ref{sub:bibliometria}; \S\ref{sub:sensibilidad} \\
14 & Selección de las fuentes de evidencia (resultados) & \S\ref{sub:res_prisma}; Figura~\ref{fig:prisma} \\
15 & Características de las fuentes & \S\ref{sub:res_estudio}; \S\ref{sub:res_complementaria} \\
16 & Evaluación crítica dentro de las fuentes & Apéndice~\ref{apx:herramienta_sesgo} (Tabla~\ref{tab:sesgo}) \\
17 & Resultados de las fuentes individuales & \S\ref{sub:res_estudio}; Tabla~\ref{tab:studies} \\
18 & Síntesis de resultados & \S\ref{sub:res_complementaria}--\ref{sub:res_mapeo} \\
19 & Resumen de la evidencia & \S\ref{sub:interpretacion} \\
20 & Limitaciones & \S\ref{sub:limitaciones} \\
21 & Conclusiones & \S\ref{sec:conclusiones} \\
22 & Financiamiento & Portada \\
\end{longtable}
\doublespacing

% ─────────────────────────────────────────────────────────────────────────────
\section{Evidencia complementaria y auditorías de sensibilidad}
\label{apx:complementaria}
% ─────────────────────────────────────────────────────────────────────────────

\subsection*{F.1 Trabajos \textit{in-domain} excluidos por criterios formales}

La Tabla~\ref{tab:complementaria} lista los 11~trabajos identificados con
la verificación léxica de \ref{sub:bibliometria} que combinan pH,
proceso azucarero y control avanzado o AA, y que no pasaron al cribado por
tipo de documento (8) o por carecer de cuartil SJR (3). Se describen
a partir de sus resúmenes y no forman parte del corpus incluido.

\singlespacing\scriptsize
\begin{longtable}{>{\raggedright\arraybackslash}p{3.0cm} >{\raggedright\arraybackslash}p{4.6cm} >{\raggedright\arraybackslash}p{2.5cm} >{\raggedright\arraybackslash}p{2.8cm} >{\raggedright\arraybackslash}p{1.9cm}}
\caption{Evidencia \textit{in-domain} excluida por criterios formales
(descripción según resumen).}\label{tab:complementaria}\\
\toprule
\textbf{Referencia} & \textbf{Proceso y método} & \textbf{Validación declarada} & \textbf{Fuente} & \textbf{Motivo} \\
\midrule
\endfirsthead
\toprule
\textbf{Referencia} & \textbf{Proceso y método} & \textbf{Validación declarada} & \textbf{Fuente} & \textbf{Motivo} \\
\midrule
\endhead
\bottomrule
\endfoot
\cite{sunori2016ga} & Neutralización de pH del clarificador; controladores ZN, IMC y LQG optimizados con GA & Simulación & \textit{Oriental J.\ Chem.} & Sin cuartil SJR \\
\cite{wang2017ph} & Clarificación en refinería; PID adaptativo difuso con diseño de \textit{hardware} y \textit{software} & Resultados experimentales (escala no precisada) & \textit{Acta Technica CSAV} & Sin cuartil SJR \\
\cite{rajalakshmi2022machine} & Clarificador; identificación y sintonía con STA, GA y PSO & No especificada & \textit{Intell.\ Autom.\ Soft Comput.} & Sin cuartil SJR (``-'') \\
\cite{sunori2017model} & Neutralizador de pH; reducción de orden de un modelo FOPDT; ZN, IMC y GA-IMC & Simulación & Acta (CICN) & Tipo documento \\
\cite{rajalakshmi2018hybrid} & Clarificador; SVR con PSO y firefly para modelado & No especificada (``problemas industriales en tiempo real'') & Acta (NPEC) & Tipo documento \\
\cite{sunori2020neurofuzzy} & Neutralización de pH; modelo neuro-difuso ANFIS & Simulación (MATLAB) & Acta (ICACCM) & Tipo documento \\
\cite{jethi2021prediction} & Neutralización de pH; MPC frente a PID & Simulación (MATLAB) & Acta (GCAT) & Tipo documento \\
\cite{kdt2021sensor} & Regulación de pH con PID-IMC; banco con celda de carga y sensor de pH & Prototipo y simulación & Acta (ICECCT) & Tipo documento \\
\cite{sunori2022digital} & Control de pH en refinería; $H_\infty$ y MPC sobre función de transferencia & Simulación (MATLAB) & Acta (ICSSIT) & Tipo documento \\
\cite{sunori2022neurofuzzy} & Control de pH en refinería; ANFIS frente a PI & Simulación (MATLAB) & Acta (ICACCS) & Tipo documento \\
\cite{sunori2022soft} & Control de pH en ingenio; PI optimizado con GA y SA & Simulación (MATLAB) & Serie \textit{LNNS} & Tipo documento \\
\end{longtable}
\doublespacing

\noindent Un registro limítrofe, \cite{vlasenko2023neural}, modela con una
red neuronal la eficiencia de la defecosaturación de jugo de remolacha
usando el pH como entrada; no se contabiliza por el criterio de
exclusión~(e).

\subsection*{F.2 Familias de expresiones regulares de la verificación léxica}

Sobre título, resumen y palabras clave en minúsculas:
\textit{pH} (\texttt{\textbackslash bph\textbackslash b}, neutralización,
alcalinización, encalado, defecación, sulfitación, carbonatación);
\textit{proceso azucarero} (\textit{sugar mill/factory/industry/process},
\textit{juice clarification}, \textit{clarifier}, \textit{raw/mixed/cane/sugar
juice}, \textit{sugar beet}, \textit{sugar refinery}, \textit{liming});
sensor virtual (\textit{soft/virtual/inferential sensor}); MPC
(\textit{model predictive control}, \textit{NMPC}, \textit{MPC});
refuerzo (\textit{reinforcement learning}, \textit{DDPG}, \textit{Q-learning},
\textit{actor-critic}); PID (\textit{PID}, \textit{proportional-integral});
difuso (\textit{fuzzy}, \textit{neuro-fuzzy}, \textit{ANFIS}, \textit{Takagi});
metaheurísticas (\textit{genetic algorithm}, \textit{particle swarm},
\textit{differential evolution}, \textit{pattern search}, \textit{state
transition algorithm}, \textit{metaheuristic}); adaptativo
(\textit{adaptive control}, \textit{gain scheduling}, \textit{self-tuning});
AA general (\textit{neural network}, \textit{machine learning}, \textit{deep
learning}, \textit{random forest}, \textit{support vector}, \textit{LSTM});
gemelo digital (\textit{digital twin}).

\subsection*{F.3 Auditoría del filtro léxico local}

De los 201~registros resueltos por el filtro L0--L2, 26 combinaban
vocabulario de pH y de proceso azucarero según un criterio más laxo
(sin reglas L0 ni requisito de AA), y 11 de ellos incluían además
vocabulario de control o estimación. Su lectura por título y resumen
mostró que tratan conservación, pardeamiento enzimático, pasteurización,
cavitación o blanqueo de jugos y bebidas, tratamiento de melazas y
residuos, filtración con membranas, dosificación de coagulantes en agua
potable o temas no relacionados (salud, fruta), y \textbf{ninguno}
aborda la regulación del pH de la clarificación o el encalado. La lista
completa (título, año, DOI y resumen) está en
\texttt{03\_filtered\_data/13.auditoria\_filtro\_local.csv}.
